\hypertarget{appendix-8bc4e8fc547c}{}
\section*{Dot products and orthogonality}
\section{Dot products and orthogonality}\label{dot-products-and-orthogonality}

\textbf{Curriculum depth:} Introduced · \textbf{First seen:} Phase6 · \textbf{Latest development:} Phase6

\subsection{Definition and mechanism}\label{definition-and-mechanism}

An inner product provides a scalar measurement of two vectors. In orthonormal real coordinates the dot is a sum of component products. Its geometric form is length product times angle cosine. Orthogonality means zero inner product; zero itself has no angle.

\((3,4,0)\cdot(2,0,0)=6\), while the cosine is \(6/(5\cdot2)=0.6\). Coordinate dotting assumes the correct metric.

\subsection{Prerequisites}\label{prerequisites}

\begin{itemize}
\tightlist
\item
  \hyperref[appendix-49933872bb51]{Norms lengths and induced distances} --- Lengths and their domains are needed before unit directions or angular measurements.
\end{itemize}

\subsection{Used by}\label{used-by}

\begin{itemize}
\item
  \href{}{Quaternion algebra and Hamilton product (01-Foundations-I, appendix-965d9260eeab)}
\item
  \hyperref[appendix-505a502c372a]{Projection rejection and scalar components} --- Builds on this mechanism in the connected geometry story.
\item
  \hyperref[appendix-3b863809f82a]{Gram matrices and coordinate metrics} --- Builds on this mechanism in the connected geometry story.
\item
  \hyperref[appendix-2391db8052b1]{Cross products and oriented area} --- Builds on this mechanism in the connected geometry story.
\item
  \hyperref[appendix-5caa058bd9c7]{Triple products and signed volume} --- Builds on this mechanism in the connected geometry story.
\item
  \hyperref[appendix-215ac7099a82]{Minkowski sums and support mappings} --- Builds on this mechanism in the connected geometry story.
\item
  \href{}{Covectors linear functionals and normal transforms (01-Foundations-I, appendix-46da295c6787)} --- Builds on this mechanism in the connected geometry story.
\item
  \hyperref[appendix-c3421e6b4f2c]{QR factorization and least squares} --- Extends this foundation into linear-systems and decomposition analysis.
\item
  \hyperref[appendix-1fad5e7f36fa]{Singular value decomposition} --- Extends this foundation into linear-systems and decomposition analysis. \#\# Recall and next depth
\end{itemize}

Derive the cosine relation and Cauchy--Schwarz; use a scale-aware orthogonality criterion. Full production proofs, target-specific integration and learner mastery require further evidence.

\subsection{Source and navigation}\label{source-and-navigation}

\href{Continuous\%20Notes/Computing\%20Foundations\%20-\%20From\%20Bits\%20to\%20C++\%20and\%20Systems\%20APIs\#The\%20dot\%20product\%20turns\%20two\%20movements\%20into\%20a\%20measurement}{Teaching narrative} · \href{}{Coverage and source gap (01-Foundations-I, appendix-fd090d982a15 § Phase 6 - Parts 501-600)} · \href{}{Evidence and limits (01-Foundations-I, appendix-a103e7846d21 § Phase 6 - Source analysis and verification)} · \hyperref[appendix-4ad47bdd269a]{Dependency Map}
